\documentclass{article}
\usepackage{graphicx} % Necessário para inserir imagens
\usepackage{amsmath}  % Para o ambiente cases
\usepackage{parskip}  % Sem recuo de parágrafo, espaço entre parágrafos
\usepackage[brazil]{babel} % Português brasileiro

\title{Clube de Matemática — Reunião 00}
\author{Herom Silva}
\date{Abril de 2026}

\begin{document}

\maketitle

\section{Introdução}
Este documento descreve a primeira reunião do Clube de Matemática e o problema que será abordado nela. Será uma reunião de uma hora e trinta minutos discutindo o problema descrito na seção Problema, que trata de uma representação matemática de Máscaras de Bits. Começaremos com uma introdução à representação de números binários.

\section{Problema}

Seja $U$ um conjunto universal finito contendo cinco variáveis:
$$U = \{x_0,\, x_1,\, x_2,\, x_3,\, x_4\}$$

Seja $S$ qualquer subconjunto de $U$ (ou seja, $S \subseteq U$). Definimos uma função característica $f(S)$, que mapeia um conjunto $S$ para uma cadeia binária de 5 bits.

A função avalia os elementos de $S$ usando uma condição lógica para gerar cada bit $b_i$ da cadeia $b_4b_3b_2b_1b_0$:
$$
b_i =
\begin{cases}
    1 & \text{se } x_i \in S \\
    0 & \text{se } x_i \notin S
\end{cases}
$$

(Nota: $b_0$ representa o bit mais à direita/menos significativo, correspondente a $x_0$.)

São dados dois subconjuntos específicos de $U$:
$$A = \{x_0,\, x_2,\, x_3\}$$
$$B = \{x_1,\, x_2,\, x_4\}$$

\begin{enumerate}
    \item Qual é $f(A)$?
    \item Qual é $f(B)$?
\end{enumerate}

\end{document}